\documentclass{article}

% if you need to pass options to natbib, use, e.g.:
% \PassOptionsToPackage{numbers, compress}{natbib}
% before loading nips_2018

% ready for submission
%\usepackage{nips_2018}

% to compile a preprint version, e.g., for submission to arXiv, add
% add the [preprint] option:
\usepackage[preprint]{nips_2018}

% to compile a camera-ready version, add the [final] option, e.g.:
% \usepackage[final]{nips_2018}

% to avoid loading the natbib package, add option nonatbib:
% \usepackage[nonatbib]{nips_2018}

\usepackage[utf8]{inputenc} % allow utf-8 input
\usepackage[T1]{fontenc}    % use 8-bit T1 fonts
\usepackage{hyperref}       % hyperlinks
\usepackage{url}            % simple URL typesetting
\usepackage{booktabs}       % professional-quality tables
\usepackage{amsfonts}       % blackboard math symbols
\usepackage{nicefrac}       % compact symbols for 1/2, etc.
\usepackage{microtype}      % microtypography

\title{Formatting instructions for NIPS 2018}
\author{
  Your Name\\
  Department of Computer Science\\
  Your University\\
  Your City, Country\\
  	exttt{your.email@university.edu}
}
  %% \And
  %% Coauthor \\
  %% Affiliation \\
  %% Address \\
  %% \texttt{email} \\
  %% \And
  %% Coauthor \\
  %% Affiliation \\
  %% Address \\
  %% \texttt{email} \\
}

\begin{document}
% \nipsfinalcopy is no longer used

\maketitle

\begin{abstract}
  The abstract paragraph should be indented \nicefrac{1}{2}~inch
  (3~picas) on both the left- and right-hand margins. Use 10~point
  type, with a vertical spacing (leading) of 11~points.  The word
  \textbf{Abstract} must be centered, bold, and in point size 12. Two
  line spaces precede the abstract. The abstract must be limited to
  one paragraph.
\end{abstract}

\section{Submission of papers to NIPS 2018}

NIPS requires electronic submissions.  The electronic submission site
is
\begin{center}
  \url{https://cmt.research.microsoft.com/NIPS2018/}
\end{center}

Please read the instructions below carefully and follow them faithfully.

\subsection{Style}

Papers to be submitted to NIPS 2018 must be prepared according to the
instructions presented here. Papers may only be up to eight pages
long, including figures. Additional pages \emph{containing only
  acknowledgments and/or cited references} are allowed. Papers that
exceed eight pages of content (ignoring references) will not be
reviewed, or in any other way considered for presentation at the
conference.

The margins in 2018 are the same as since 2007, which allow for
$\sim$$15\%$ more words in the paper compared to earlier years.

Authors are required to use the NIPS \LaTeX{} style files obtainable
at the NIPS website as indicated below. Please make sure you use the
current files and not previous versions. Tweaking the style files may
be grounds for rejection.

\subsection{Retrieval of style files}

The style files for NIPS and other conference information are
available on the World Wide Web at
\begin{center}
  \url{http://www.nips.cc/}
\end{center}
The file \verb+nips_2018.pdf+ contains these instructions and
illustrates the various formatting requirements your NIPS paper must
satisfy.

The only supported style file for NIPS 2018 is \verb+nips_2018.sty+,
rewritten for \LaTeXe{}.  \textbf{Previous style files for \LaTeX{}
  2.09, Microsoft Word, and RTF are no longer supported!}

The \LaTeX{} style file contains three optional arguments: \verb+final+,
which creates a camera-ready copy, \verb+preprint+, which creates a
preprint for submission to, e.g., arXiv, and \verb+nonatbib+, which will
not load the \verb+natbib+ package for you in case of package clash.

\paragraph{New preprint option for 2018}
If you wish to post a preprint of your work online, e.g., on arXiv,
using the NIPS style, please use the \verb+preprint+ option. This will
create a nonanonymized version of your work with the text
``Preprint. Work in progress.''  in the footer. This version may be
distributed as you see fit. Please \textbf{do not} use the
\verb+final+ option, which should \textbf{only} be used for papers
accepted to NIPS.

At submission time, please omit the \verb+final+ and \verb+preprint+
options. This will anonymize your submission and add line numbers to aid
review. Please do \emph{not} refer to these line numbers in your paper
as they will be removed during generation of camera-ready copies.

The file \verb+nips_2018.tex+ may be used as a ``shell'' for writing
your paper. All you have to do is replace the author, title, abstract,
and text of the paper with your own.

The formatting instructions contained in these style files are
summarized in Sections \ref{gen_inst}, \ref{headings}, and
\ref{others} below.

\section{General formatting instructions}
\label{gen_inst}

The text must be confined within a rectangle 5.5~inches (33~picas)
wide and 9~inches (54~picas) long. The left margin is 1.5~inch
(9~picas).  Use 10~point type with a vertical spacing (leading) of
11~points.  Times New Roman is the preferred typeface throughout, and
will be selected for you by default.  Paragraphs are separated by
\nicefrac{1}{2}~line space (5.5 points), with no indentation.

The paper title should be 17~point, initial caps/lower case, bold,
centered between two horizontal rules. The top rule should be 4~points
thick and the bottom rule should be 1~point thick. Allow
\nicefrac{1}{4}~inch space above and below the title to rules. All
pages should start at 1~inch (6~picas) from the top of the page.

For the final version, authors' names are set in boldface, and each
name is centered above the corresponding address. The lead author's
name is to be listed first (left-most), and the co-authors' names (if
different address) are set to follow. If there is only one co-author,
list both author and co-author side by side.

Please pay special attention to the instructions in Section \ref{others}
regarding figures, tables, acknowledgments, and references.

\section{Headings: first level}
\label{headings}

All headings should be lower case (except for first word and proper
nouns), flush left, and bold.

First-level headings should be in 12-point type.

\subsection{Headings: second level}

Second-level headings should be in 10-point type.

\subsubsection{Headings: third level}

Third-level headings should be in 10-point type.

\paragraph{Paragraphs}

There is also a \verb+\paragraph+ command available, which sets the
heading in bold, flush left, and inline with the text, with the
heading followed by 1\,em of space.

\section{Citations, figures, tables, references}
\label{others}

These instructions apply to everyone.

\subsection{Citations within the text}

The \verb+natbib+ package will be loaded for you by default.
Citations may be author/year or numeric, as long as you maintain
internal consistency.  As to the format of the references themselves,
any style is acceptable as long as it is used consistently.

The documentation for \verb+natbib+ may be found at
\begin{center}
  \url{http://mirrors.ctan.org/macros/latex/contrib/natbib/natnotes.pdf}
\end{center}
Of note is the command \verb+\citet+, which produces citations
appropriate for use in inline text.  For example,
\begin{verbatim}
   \citet{hasselmo} investigated\dots
\end{verbatim}
produces
\begin{quote}
  Hasselmo, et al.\ (1995) investigated\dots
\end{quote}

If you wish to load the \verb+natbib+ package with options, you may
add the following before loading the \verb+nips_2018+ package:
\begin{verbatim}
   \PassOptionsToPackage{options}{natbib}
\end{verbatim}

If \verb+natbib+ clashes with another package you load, you can add
the optional argument \verb+nonatbib+ when loading the style file:
\begin{verbatim}
   \usepackage[nonatbib]{nips_2018}
\end{verbatim}

As submission is double blind, refer to your own published work in the
third person. That is, use ``In the previous work of Jones et
al.\ [4],'' not ``In our previous work [4].'' If you cite your other
papers that are not widely available (e.g., a journal paper under
review), use anonymous author names in the citation, e.g., an author
of the form ``A.\ Anonymous.''

\subsection{Footnotes}

Footnotes should be used sparingly.  If you do require a footnote,
indicate footnotes with a number\footnote{Sample of the first
  footnote.} in the text. Place the footnotes at the bottom of the
page on which they appear.  Precede the footnote with a horizontal
rule of 2~inches (12~picas).

Note that footnotes are properly typeset \emph{after} punctuation
marks.\footnote{As in this example.}

\subsection{Figures}

\begin{figure}
  \centering
  \fbox{\rule[-.5cm]{0cm}{4cm} \rule[-.5cm]{4cm}{0cm}}
  \caption{Sample figure caption.}
\end{figure}

All artwork must be neat, clean, and legible. Lines should be dark
enough for purposes of reproduction. The figure number and caption
always appear after the figure. Place one line space before the figure
caption and one line space after the figure. The figure caption should
be lower case (except for first word and proper nouns); figures are
numbered consecutively.

You may use color figures.  However, it is best for the figure
captions and the paper body to be legible if the paper is printed in
either black/white or in color.

\subsection{Tables}

All tables must be centered, neat, clean and legible.  The table
number and title always appear before the table.  See
Table~\ref{sample-table}.

Place one line space before the table title, one line space after the
table title, and one line space after the table. The table title must
be lower case (except for first word and proper nouns); tables are
numbered consecutively.

Note that publication-quality tables \emph{do not contain vertical
  rules.} We strongly suggest the use of the \verb+booktabs+ package,
which allows for typesetting high-quality, professional tables:
\begin{center}
  \url{https://www.ctan.org/pkg/booktabs}
\end{center}
This package was used to typeset Table~\ref{sample-table}.

\begin{table}
  \caption{Sample table title}
  \label{sample-table}
  \centering
  \begin{tabular}{lll}
    \toprule
    \multicolumn{2}{c}{Part}                   \\
    \cmidrule(r){1-2}
    Name     & Description     & Size ($\mu$m) \\
    \midrule
    Dendrite & Input terminal  & $\sim$100     \\
    Axon     & Output terminal & $\sim$10      \\
    Soma     & Cell body       & up to $10^6$  \\
    \bottomrule
  \end{tabular}
\end{table}

\section{Mathematical Notation and Equations}

\subsection{Basic Mathematical Formatting}

Mathematical expressions should be formatted using \LaTeX{}'s math mode.
Inline mathematics should be enclosed in single dollar signs (\verb+$+),
while displayed equations should be enclosed in double dollar signs
(\verb+$$+) or use the \verb+\[+\ and \verb+\]+ delimiters.

For example, the quadratic formula $x = \frac{-b \pm \sqrt{b^2 - 4ac}}{2a}$
can be displayed as:
\[
x = \frac{-b \pm \sqrt{b^2 - 4ac}}{2a}
\]

\subsection{Commonly Used Symbols}

Here are some commonly used mathematical symbols in NIPS papers:

\begin{itemize}
\item Sets and spaces: $\mathbb{R}$, $\mathbb{N}$, $\mathbb{C}$, $\mathcal{X}$, $\mathcal{Y}$
\item Operators: $\nabla$, $\partial$, $\Delta$, $\div$, $\curl$
\item Probability: $\mathbb{P}$, $\mathbb{E}$, $\mathcal{N}$, $\mu$, $\sigma^2$
\item Information theory: $H(X)$, $I(X;Y)$, $D_{KL}(p||q)$
\item Linear algebra: $\mathbf{x}$, $\mathbf{W}$, $\mathbf{b}$, $\mathbf{A}^\top$
\end{itemize}

\subsection{Theorem Environments}

Theorems, lemmas, propositions, and corollaries should be formatted using
the appropriate environments:

\begin{theorem}
Let $f: \mathbb{R}^n \to \mathbb{R}$ be a differentiable function. Then
the gradient $\nabla f(\mathbf{x})$ points in the direction of steepest
ascent of $f$ at $\mathbf{x}$.
\end{theorem}

\begin{proof}
The proof follows directly from the definition of the directional
derivative and the Cauchy-Schwarz inequality.
\end{proof}

\begin{lemma}
If $\mathbf{A}$ is a symmetric positive definite matrix, then there
exists a unique Cholesky decomposition $\mathbf{A} = \mathbf{L}\mathbf{L}^\top$
where $\mathbf{L}$ is lower triangular with positive diagonal elements.
\end{lemma}

\subsection{Algorithm Formatting}

Algorithms should be formatted using the algorithmic package or similar
environments. Here's an example:

\begin{algorithmic}[1]
\Procedure{GradientDescent}{$\mathbf{x}_0, \eta, T$}
    \State $\mathbf{x} \gets \mathbf{x}_0$
    \For{$t = 1$ \textbf{to} $T$}
        \State $\mathbf{x} \gets \mathbf{x} - \eta \nabla f(\mathbf{x})$
    \EndFor
    \State \textbf{return} $\mathbf{x}$
\EndProcedure
\end{algorithmic}

\section{Machine Learning Concepts}

\subsection{Supervised Learning}

In supervised learning, we aim to learn a mapping from input features
$\mathbf{x} \in \mathcal{X}$ to output labels $y \in \mathcal{Y}$ using
a training dataset $\mathcal{D} = \{(\mathbf{x}_i, y_i)\}_{i=1}^n$.

The general supervised learning objective is to minimize the expected
risk:
\[
\mathcal{R}(h) = \mathbb{E}_{(\mathbf{x},y) \sim \mathcal{P}}[\ell(h(\mathbf{x}), y)]
\]
where $\ell$ is a loss function and $h$ is the hypothesis.

\subsection{Unsupervised Learning}

Unsupervised learning deals with data $\mathbf{x} \in \mathcal{X}$ without
corresponding labels. Common approaches include:

\begin{enumerate}
\item \textbf{Clustering}: Partition data into groups with similar properties
\item \textbf{Dimensionality Reduction}: Find lower-dimensional representations
\item \textbf{Density Estimation}: Model the underlying data distribution $p(\mathbf{x})$
\item \textbf{Generative Modeling}: Learn to generate new samples from $p(\mathbf{x})$
\end{enumerate}

\subsection{Deep Learning Architectures}

\subsubsection{Convolutional Neural Networks}

Convolutional Neural Networks (CNNs) are particularly effective for
processing grid-structured data such as images. A convolutional layer
computes:
\[
\mathbf{y}_{i,j,k} = \sum_{m=1}^M \sum_{p=0}^{P-1} \sum_{q=0}^{Q-1}
\mathbf{x}_{i+p,j+q,m} \mathbf{w}_{p,q,m,k} + b_k
\]

\subsubsection{Recurrent Neural Networks}

Recurrent Neural Networks (RNNs) are designed for sequential data.
The basic RNN update is:
\[
\mathbf{h}_t = \tanh(\mathbf{W}_{xh} \mathbf{x}_t + \mathbf{W}_{hh} \mathbf{h}_{t-1} + \mathbf{b}_h)
\]
\[
\mathbf{y}_t = \mathbf{W}_{hy} \mathbf{h}_t + \mathbf{b}_y
\]

\subsubsection{Attention Mechanisms}

Attention mechanisms allow models to focus on relevant parts of the input:
\[
\alpha_{i,j} = \frac{\exp(e_{i,j})}{\sum_{k=1}^L \exp(e_{i,k})}
\]
\[
\mathbf{c}_i = \sum_{j=1}^L \alpha_{i,j} \mathbf{h}_j
\]

\subsection{Optimization Algorithms}

\subsubsection{Gradient Descent Variants}

\begin{itemize}
\item \textbf{Batch Gradient Descent}: Uses entire dataset
\[
\mathbf{w} \leftarrow \mathbf{w} - \eta \frac{1}{n} \sum_{i=1}^n \nabla_\mathbf{w} \ell(\mathbf{w}; \mathbf{x}_i, y_i)
\]

\item \textbf{Stochastic Gradient Descent}: Uses single example
\[
\mathbf{w} \leftarrow \mathbf{w} - \eta \nabla_\mathbf{w} \ell(\mathbf{w}; \mathbf{x}_i, y_i)
\]

\item \textbf{Mini-batch Gradient Descent}: Uses small batches
\[
\mathbf{w} \leftarrow \mathbf{w} - \eta \frac{1}{B} \sum_{i=1}^B \nabla_\mathbf{w} \ell(\mathbf{w}; \mathbf{x}_{i+b}, y_{i+b})
\]
\end{itemize}

\subsubsection{Adaptive Optimization}

Adam combines momentum and RMSProp:
\[
\mathbf{m}_t = \beta_1 \mathbf{m}_{t-1} + (1 - \beta_1) \mathbf{g}_t
\]
\[
\mathbf{v}_t = \beta_2 \mathbf{v}_{t-1} + (1 - \beta_2) \mathbf{g}_t^2
\]
\[
\hat{\mathbf{m}}_t = \frac{\mathbf{m}_t}{1 - \beta_1^t}, \quad
\hat{\mathbf{v}}_t = \frac{\mathbf{v}_t}{1 - \beta_2^t}
\]
\[
\mathbf{w}_{t+1} = \mathbf{w}_t - \eta \frac{\hat{\mathbf{m}}_t}{\sqrt{\hat{\mathbf{v}}_t} + \epsilon}
\]

\section{Statistical Methods}

\subsection{Bayesian Inference}

Bayesian inference provides a framework for updating beliefs about
parameters given data. The posterior distribution is:
\[
p(\theta|\mathcal{D}) = \frac{p(\mathcal{D}|\theta) p(\theta)}{p(\mathcal{D})}
\]

\subsubsection{Conjugate Priors}

For certain likelihood-prior combinations, the posterior has the same
functional form as the prior:

\begin{itemize}
\item Beta-Binomial: $p(\theta) = \Beta(\alpha, \beta)$
\item Normal-Normal: $p(\theta) = \mathcal{N}(\mu_0, \sigma_0^2)$
\item Gamma-Poisson: $p(\theta) = \Gamma(\alpha, \beta)$
\end{itemize}

\subsubsection{Markov Chain Monte Carlo}

MCMC methods generate samples from complex distributions. The Metropolis-Hastings
algorithm:

\begin{enumerate}
\item Sample $\theta' \sim q(\theta'|\theta^{(t)})$
\item Compute acceptance probability:
\[
\alpha = \min\left(1, \frac{p(\theta'|\mathcal{D}) q(\theta^{(t)}|\theta')}{p(\theta^{(t)}|\mathcal{D}) q(\theta'|\theta^{(t)})}\right)
\]
\item Set $\theta^{(t+1)} = \theta'$ with probability $\alpha$, otherwise $\theta^{(t+1)} = \theta^{(t)}$
\end{enumerate}

\subsection{Hypothesis Testing}

\subsubsection{p-values and Confidence Intervals}

A p-value is the probability of observing data at least as extreme as
the observed data, assuming the null hypothesis is true:
\[
p\text{-value} = \mathbb{P}(T \geq t|\mathcal{H}_0)
\]

For a 95\% confidence interval:
\[
\hat{\theta} \pm 1.96 \cdot \SE(\hat{\theta})
\]

\subsubsection{Multiple Testing Correction}

When performing multiple hypothesis tests, we need to control the
family-wise error rate or false discovery rate.

Bonferroni correction: $\alpha' = \alpha / m$

\subsubsection{Power Analysis}

The power of a test is the probability of correctly rejecting a false
null hypothesis:
\[
\text{Power} = \mathbb{P}(\text{reject } \mathcal{H}_0 | \mathcal{H}_1 \text{ is true})
\]

\section{Information Theory}

\subsection{Entropy and Mutual Information}

The entropy of a discrete random variable $X$ is:
\[
H(X) = -\sum_{x} p(x) \log p(x)
\]

The mutual information between $X$ and $Y$ is:
\[
I(X;Y) = H(X) + H(Y) - H(X,Y)
\]

\subsection{Kullback-Leibler Divergence}

The KL divergence measures the difference between two distributions:
\[
D_{KL}(p||q) = \sum_x p(x) \log \frac{p(x)}{q(x)}
\]

Properties of KL divergence:
\begin{itemize}
\item $D_{KL}(p||q) \geq 0$
\item $D_{KL}(p||q) = 0$ iff $p = q$
\item Not symmetric: $D_{KL}(p||q) \neq D_{KL}(q||p)$
\end{itemize}

\subsection{Cross-Entropy Loss}

Cross-entropy loss is commonly used for classification:
\[
\ell(\hat{y}, y) = -\sum_{k=1}^K y_k \log \hat{y}_k
\]

For binary classification, this reduces to:
\[
\ell(\hat{y}, y) = -[y \log \hat{y} + (1-y) \log (1-\hat{y})]
\]

\section{Linear Algebra Review}

\subsection{Vector and Matrix Operations}

\subsubsection{Inner Products and Norms}

The inner product between vectors $\mathbf{u}, \mathbf{v} \in \mathbb{R}^n$:
\[
\langle \mathbf{u}, \mathbf{v} \rangle = \mathbf{u}^\top \mathbf{v} = \sum_{i=1}^n u_i v_i
\]

The $\ell_2$ norm:
\[
\|\mathbf{v}\|_2 = \sqrt{\langle \mathbf{v}, \mathbf{v} \rangle} = \sqrt{\sum_{i=1}^n v_i^2}
\]

\subsubsection{Matrix Decompositions}

Singular Value Decomposition (SVD):
\[
\mathbf{A} = \mathbf{U} \mathbf{\Sigma} \mathbf{V}^\top
\]
where $\mathbf{U}$ and $\mathbf{V}$ are orthogonal matrices, and
$\mathbf{\Sigma}$ is diagonal with singular values.

Eigenvalue decomposition for symmetric matrices:
\[
\mathbf{A} = \mathbf{Q} \mathbf{\Lambda} \mathbf{Q}^\top
\]

\subsubsection{Matrix Derivatives}

Useful derivatives for optimization:

\begin{align*}
\frac{\partial}{\partial \mathbf{x}} (\mathbf{x}^\top \mathbf{A} \mathbf{x}) &= 2\mathbf{A} \mathbf{x} \\
\frac{\partial}{\partial \mathbf{A}} \|\mathbf{A} \mathbf{x} - \mathbf{b}\|_2^2 &= 2(\mathbf{A} \mathbf{x} - \mathbf{b}) \mathbf{x}^\top \\
\frac{\partial}{\partial \mathbf{W}} \log \det(\mathbf{W}) &= (\mathbf{W}^{-1})^\top
\end{align*}

\section{Probability Distributions}

\subsection{Common Distributions}

\subsubsection{Gaussian Distribution}

The multivariate Gaussian distribution:
\[
\mathcal{N}(\mathbf{x}; \boldsymbol{\mu}, \boldsymbol{\Sigma}) =
\frac{1}{(2\pi)^{d/2} |\boldsymbol{\Sigma}|^{1/2}}
\exp\left( -\frac{1}{2} (\mathbf{x} - \boldsymbol{\mu})^\top \boldsymbol{\Sigma}^{-1} (\mathbf{x} - \boldsymbol{\mu}) \right)
\]

\subsubsection{Bernoulli Distribution}

For binary random variables:
\[
\text{Bern}(x; p) = p^x (1-p)^{1-x}
\]

\subsubsection{Categorical Distribution}

For discrete random variables with $K$ categories:
\[
\text{Cat}(\mathbf{x}; \mathbf{p}) = \prod_{k=1}^K p_k^{x_k}
\]

\subsection{Convergence Concepts}

\subsubsection{Laws of Large Numbers}

Weak Law of Large Numbers: For i.i.d. random variables with finite mean $\mu$:
\[
\frac{1}{n} \sum_{i=1}^n X_i \xrightarrow{P} \mu
\]

Strong Law of Large Numbers:
\[
\frac{1}{n} \sum_{i=1}^n X_i \xrightarrow{a.s.} \mu
\]

\subsubsection{Central Limit Theorem}

For i.i.d. random variables with finite mean $\mu$ and variance $\sigma^2$:
\[
\sqrt{n} \left( \frac{1}{n} \sum_{i=1}^n X_i - \mu \right) \xrightarrow{D} \mathcal{N}(0, \sigma^2)
\]

\subsection{Concentration Inequalities}

\subsubsection{Chebyshev's Inequality}

For any random variable $X$ with finite variance:
\[
\mathbb{P}(|X - \mathbb{E}[X]| \geq t) \leq \frac{\Var(X)}{t^2}
\]

\subsubsection{Hoeffding's Inequality}

For bounded independent random variables $X_i \in [a_i, b_i]$:
\[
\mathbb{P}\left( \left| \frac{1}{n} \sum_{i=1}^n X_i - \mathbb{E}\left[\frac{1}{n} \sum_{i=1}^n X_i\right] \right| \geq t \right) \leq 2\exp\left( -\frac{2n^2 t^2}{\sum_{i=1}^n (b_i - a_i)^2} \right)
\]

\section{Computational Complexity}

\subsection{Time Complexity}

Common time complexity classes:
\begin{itemize}
\item $\mathcal{O}(1)$: Constant time
\item $\mathcal{O}(\log n)$: Logarithmic time
\item $\mathcal{O}(n)$: Linear time
\item $\mathcal{O}(n \log n)$: Linearithmic time
\item $\mathcal{O}(n^2)$: Quadratic time
\item $\mathcal{O}(2^n)$: Exponential time
\end{itemize}

\subsection{Space Complexity}

Space complexity measures the amount of memory required by an algorithm.

\subsubsection{In-place Algorithms}

Algorithms that use $\mathcal{O}(1)$ additional space beyond the input.

\subsubsection{Memoization}

Dynamic programming technique that stores results of expensive function calls:
\[
T(n) = \begin{cases}
\Theta(1) & \text{if } n \in \mathcal{B} \\
\Theta(n) + \sum_{i=1}^k T(n_i) & \text{otherwise}
\end{cases}
\]

\section{Advanced Topics}

\subsection{Variational Inference}

Variational inference approximates complex posterior distributions with
simpler distributions from a family $\mathcal{Q}$:
\[
p(\theta|\mathcal{D}) \approx q(\theta) = \arg\min_{q \in \mathcal{Q}} D_{KL}(q(\theta)||p(\theta|\mathcal{D}))
\]

The Evidence Lower Bound (ELBO):
\[
\log p(\mathcal{D}) \geq \mathbb{E}_{q(\theta)}[\log p(\mathcal{D},\theta)] - \mathbb{E}_{q(\theta)}[\log q(\theta)]
\]

\subsection{Generative Adversarial Networks}

GANs consist of a generator $G$ and discriminator $D$:
\[
\min_G \max_D \mathbb{E}_{x \sim p_{data}}[\log D(x)] + \mathbb{E}_{z \sim p_z}[\log (1 - D(G(z)))]
\]

The optimal discriminator is:
\[
D^*(x) = \frac{p_{data}(x)}{p_{data}(x) + p_g(x)}
\]

\subsection{Reinforcement Learning}

The Bellman equation for value functions:
\[
V^\pi(s) = \mathbb{E}_\pi [r(s,a) + \gamma V^\pi(s')]
\]

Q-learning update:
\[
Q(s,a) \leftarrow Q(s,a) + \alpha \left( r + \gamma \max_{a'} Q(s',a') - Q(s,a) \right)
\]

\section{Additional Examples}

\subsection{More Figure Examples}

\begin{figure}[h]
\centering
\subfigure[Training Loss]{
\begin{minipage}[b]{0.45\textwidth}
\centering
\fbox{\rule[-.5cm]{0cm}{3cm} \rule[-.5cm]{3cm}{0cm}}
\caption{Training loss over epochs}
\label{fig:train_loss}
\end{minipage}
}
\subfigure[Validation Accuracy]{
\begin{minipage}[b]{0.45\textwidth}
\centering
\fbox{\rule[-.5cm]{0cm}{3cm} \rule[-.5cm]{3cm}{0cm}}
\caption{Validation accuracy over epochs}
\label{fig:val_acc}
\end{minipage}
}
\caption{Training curves for a neural network model}
\label{fig:training_curves}
\end{figure}

\begin{figure}[h]
\centering
\fbox{\rule[-.5cm]{0cm}{4cm} \rule[-.5cm]{4cm}{0cm}}
\caption{Architecture diagram of a convolutional neural network}
\label{fig:cnn_arch}
\end{figure}

\subsection{More Table Examples}

\begin{table}[h]
\caption{Comparison of optimization algorithms}
\label{tab:opt_comparison}
\centering
\begin{tabular}{@{}llll@{}}
\toprule
Algorithm & Convergence Rate & Memory Usage & Adaptive Learning Rate \\
\midrule
SGD & $\mathcal{O}(1/\sqrt{T})$ & $\mathcal{O}(d)$ & No \\
Momentum & $\mathcal{O}(1/\sqrt{T})$ & $\mathcal{O}(d)$ & No \\
Adagrad & $\mathcal{O}(1/\sqrt{T})$ & $\mathcal{O}(d)$ & Yes \\
RMSProp & $\mathcal{O}(1/\sqrt{T})$ & $\mathcal{O}(d)$ & Yes \\
Adam & $\mathcal{O}(1/\sqrt{T})$ & $\mathcal{O}(d)$ & Yes \\
\bottomrule
\end{tabular}
\end{table}

\begin{table}[h]
\caption{Hyperparameter search results}
\label{tab:hyperparams}
\centering
\begin{tabular}{@{}llll@{}}
\toprule
Learning Rate & Batch Size & Hidden Units & Validation Accuracy \\
\midrule
0.001 & 32 & 128 & 85.2\% \\
0.001 & 64 & 128 & 86.1\% \\
0.001 & 32 & 256 & 87.3\% \\
0.001 & 64 & 256 & 88.7\% \\
0.01 & 32 & 128 & 82.1\% \\
0.01 & 64 & 128 & 83.9\% \\
\bottomrule
\end{tabular}
\end{table}

\section{Code Examples}

\subsection{Python Code for Neural Networks}

\begin{verbatim}
import torch
import torch.nn as nn
import torch.nn.functional as F

class SimpleCNN(nn.Module):
    def __init__(self, num_classes=10):
        super(SimpleCNN, self).__init__()
        self.conv1 = nn.Conv2d(3, 32, kernel_size=3, padding=1)
        self.conv2 = nn.Conv2d(32, 64, kernel_size=3, padding=1)
        self.fc1 = nn.Linear(64 * 8 * 8, 512)
        self.fc2 = nn.Linear(512, num_classes)

    def forward(self, x):
        x = F.relu(self.conv1(x))
        x = F.max_pool2d(x, 2)
        x = F.relu(self.conv2(x))
        x = F.max_pool2d(x, 2)
        x = x.view(-1, 64 * 8 * 8)
        x = F.relu(self.fc1(x))
        x = F.dropout(x, training=self.training)
        x = self.fc2(x)
        return x
\end{verbatim}

\subsection{Mathematical Derivations}

\subsubsection{Backpropagation}

The backpropagation algorithm computes gradients through the chain rule:

For a loss function $\ell$ and network $f(\mathbf{x}; \theta)$, the gradient
with respect to parameters $\theta$ is:
\[
\frac{\partial \ell}{\partial \theta} = \frac{\partial \ell}{\partial f} \cdot \frac{\partial f}{\partial \theta}
\]

For a multi-layer network:
\[
\frac{\partial \ell}{\partial \mathbf{W}^{(L)}} = \delta^{(L)} (\mathbf{h}^{(L-1)})^\top
\]
\[
\delta^{(l)} = (\mathbf{W}^{(l+1)})^\top \delta^{(l+1)} \odot \sigma'(\mathbf{z}^{(l)})
\]

\subsubsection{Batch Normalization}

Batch normalization normalizes layer inputs to have zero mean and unit variance:
\[
\hat{\mathbf{x}}^{(k)} = \frac{\mathbf{x}^{(k)} - \mathbb{E}[\mathbf{x}^{(k)}]}{\sqrt{\Var(\mathbf{x}^{(k)}) + \epsilon}}
\]
\[
\mathbf{y}^{(k)} = \gamma^{(k)} \hat{\mathbf{x}}^{(k)} + \beta^{(k)}
\]

During inference, running statistics are used:
\[
\hat{\mathbf{x}}^{(k)} = \frac{\mathbf{x}^{(k)} - \hat{\mu}^{(k)}}{\sqrt{\hat{\sigma}^{2(k)} + \epsilon}}
\]

\section{Experimental Methodology}

\subsection{Cross-Validation}

\subsubsection{k-Fold Cross-Validation}

In k-fold cross-validation, the dataset is divided into k equal parts.
For each fold:

\begin{enumerate}
\item Use fold i as test set
\item Use remaining k-1 folds as training set
\item Train model and evaluate on test set
\item Repeat for all k folds
\item Average performance across all folds
\end{enumerate}

\subsubsection{Stratified k-Fold}

For imbalanced datasets, stratified k-fold ensures each fold has the
same proportion of each class as the original dataset.

\subsection{Baseline Methods}

Always compare against appropriate baselines:

\begin{itemize}
\item \textbf{Random baseline}: Random predictions
\item \textbf{Majority class}: Always predict most frequent class
\item \textbf{Simple models}: Logistic regression, decision trees
\item \textbf{Domain-specific}: State-of-the-art methods for the task
\end{itemize}

\subsection{Evaluation Metrics}

\subsubsection{Classification Metrics}

For binary classification:
\begin{align*}
\text{Accuracy} &= \frac{TP + TN}{TP + TN + FP + FN} \\
\text{Precision} &= \frac{TP}{TP + FP} \\
\text{Recall} &= \frac{TP}{TP + FN} \\
F_1 &= 2 \cdot \frac{\text{Precision} \cdot \text{Recall}}{\text{Precision} + \text{Recall}}
\end{align*}

\subsubsection{Regression Metrics}

Common regression metrics:
\begin{align*}
\text{MSE} &= \frac{1}{n} \sum_{i=1}^n (y_i - \hat{y}_i)^2 \\
\text{MAE} &= \frac{1}{n} \sum_{i=1}^n |y_i - \hat{y}_i| \\
R^2 &= 1 - \frac{\sum (y_i - \hat{y}_i)^2}{\sum (y_i - \bar{y})^2}
\end{align*}

\subsection{Statistical Significance Testing}

\subsubsection{Paired t-test}

For comparing two models on the same test set:
\[
t = \frac{\bar{d}}{s_d / \sqrt{n}}
\]
where $d_i = \text{performance}_i^A - \text{performance}_i^B$

\subsubsection{McNemar's Test}

For comparing classification accuracy on the same instances:
\[
\chi^2 = \frac{(|n_{01} - n_{10}| - 1)^2}{n_{01} + n_{10}}
\]

\section{Practical Considerations}

\subsection{Data Preprocessing}

\subsubsection{Feature Scaling}

Standardization (z-score normalization):
\[
\mathbf{x}' = \frac{\mathbf{x} - \mu}{\sigma}
\]

Min-max normalization:
\[
\mathbf{x}' = \frac{\mathbf{x} - \min(\mathbf{x})}{\max(\mathbf{x}) - \min(\mathbf{x})}
\]

\subsubsection{Handling Missing Data}

Common strategies:
\begin{itemize}
\item \textbf{Deletion}: Remove instances with missing values
\item \textbf{Imputation}: Fill missing values with mean, median, or mode
\item \textbf{Model-based}: Use algorithms that handle missing data
\end{itemize}

\subsection{Regularization Techniques}

\subsubsection{L2 Regularization}

Adds penalty term to loss function:
\[
\ell'(\theta) = \ell(\theta) + \lambda \|\theta\|_2^2
\]

The gradient becomes:
\[
\frac{\partial \ell'}{\partial \theta} = \frac{\partial \ell}{\partial \theta} + 2\lambda \theta
\]

\subsubsection{Dropout}

Randomly sets activations to zero during training:
\[
\mathbf{h}' = \mathbf{h} \odot \mathbf{m}
\]
where $\mathbf{m} \sim \text{Bernoulli}(p)$

\subsection{Hyperparameter Optimization}

\subsubsection{Grid Search}

Exhaustive search over hyperparameter combinations:
\[
\mathcal{H} = \{\eta, \lambda, d\} \times \{0.001, 0.01, 0.1\} \times \{0.0001, 0.001, 0.01\} \times \{64, 128, 256\}
\]

\subsubsection{Random Search}

Sample hyperparameters randomly from distributions:
\[
\eta \sim \log\mathcal{U}(10^{-4}, 10^{-1})
\]
\[
\lambda \sim \log\mathcal{U}(10^{-5}, 10^{-3})
\]

\subsubsection{Bayesian Optimization}

Use Gaussian processes to model the objective function and select
next hyperparameters to evaluate.

\section{Ethical Considerations}

\subsection{Bias and Fairness}

\subsubsection{Dataset Bias}

Models can inherit biases from training data:
\begin{itemize}
\item \textbf{Selection bias}: Non-representative training data
\item \textbf{Label bias}: Incorrect or biased annotations
\item \textbf{Measurement bias}: Systematic measurement errors
\end{itemize}

\subsubsection{Fairness Metrics}

Demographic parity:
\[
\mathbb{P}(\hat{y} = 1 | A = a) = \mathbb{P}(\hat{y} = 1 | A = b)
\]

Equal opportunity:
\[
\mathbb{P}(\hat{y} = 1 | A = a, Y = 1) = \mathbb{P}(\hat{y} = 1 | A = b, Y = 1)
\]

\subsection{Privacy and Security}

\subsubsection{Differential Privacy}

A mechanism $\mathcal{M}$ is $(\epsilon, \delta)$-differentially private if:
\[
\mathbb{P}(\mathcal{M}(D) \in S) \leq e^\epsilon \mathbb{P}(\mathcal{M}(D') \in S) + \delta
\]
for all neighboring datasets $D, D'$ and all $S \subseteq \text{Range}(\mathcal{M})$

\subsubsection{Adversarial Examples}

Small perturbations that can fool neural networks:
\[
\mathbf{x}' = \mathbf{x} + \epsilon \cdot \sign(\nabla_\mathbf{x} \ell(f(\mathbf{x}), y))
\]

\subsection{Interpretability}

\subsubsection{Feature Importance}

Permutation feature importance:
\begin{enumerate}
\item Train model on original data
\item For each feature, permute its values across instances
\item Measure decrease in model performance
\item Feature importance = original performance - permuted performance
\end{enumerate}

\subsubsection{SHAP Values}

Shapley Additive Explanations assign each feature an importance value:
\[
\phi_i = \sum_{S \subseteq N \setminus \{i\}} \frac{|S|! (|N| - |S| - 1)!}{|N|!}
\left[ f(S \cup \{i\}) - f(S) \right]
\]

\section{Advanced LaTeX Features}

\subsection{Custom Commands}

You can define custom commands for frequently used expressions:

\begin{verbatim}
\newcommand{\R}{\mathbb{R}}
\newcommand{\E}{\mathbb{E}}
\newcommand{\Var}{\mathrm{Var}}
\newcommand{\Cov}{\mathrm{Cov}}
\end{verbatim}

\subsection{Bibliography Management}

Use BibTeX for managing references. Example bib entry:

\begin{verbatim}
@article{smith2020deep,
  title={Deep Learning for Natural Language Processing},
  author={Smith, John and Doe, Jane},
  journal={Journal of Machine Learning},
  volume={15},
  number={2},
  pages={123--145},
  year={2020},
  publisher={ML Press}
}
\end{verbatim}

\subsection{Acronyms and Glossaries}

Define acronyms for technical terms:

\begin{verbatim}
\usepackage{acronym}
\begin{acronym}
  \acro{ML}{Machine Learning}
  \acro{DL}{Deep Learning}
  \acro{NN}{Neural Network}
  \acro{CNN}{Convolutional Neural Network}
  \acro{RNN}{Recurrent Neural Network}
\end{acronym}
\end{verbatim}

Use in text: \ac{ML}, \acf{DL} (first use), \acp{CNN} (plural).

\section{Final instructions}

Do not change any aspects of the formatting parameters in the style
files.  In particular, do not modify the width or length of the
rectangle the text should fit into, and do not change font sizes
(except perhaps in the \textbf{References} section; see below). Please
note that pages should be numbered.

\section{Preparing PDF files}

Please prepare submission files with paper size ``US Letter,'' and
not, for example, ``A4.''

Fonts were the main cause of problems in the past years. Your PDF file
must only contain Type 1 or Embedded TrueType fonts. Here are a few
instructions to achieve this.

\begin{itemize}

\item You should directly generate PDF files using \verb+pdflatex+.

\item You can check which fonts a PDF files uses.  In Acrobat Reader,
  select the menu Files$>$Document Properties$>$Fonts and select Show
  All Fonts. You can also use the program \verb+pdffonts+ which comes
  with \verb+xpdf+ and is available out-of-the-box on most Linux
  machines.

\item The IEEE has recommendations for generating PDF files whose
  fonts are also acceptable for NIPS. Please see
  \url{http://www.emfield.org/icuwb2010/downloads/IEEE-PDF-SpecV32.pdf}

\item \verb+xfig+ "patterned" shapes are implemented with bitmap
  fonts.  Use "solid" shapes instead.

\item The \verb+\bbold+ package almost always uses bitmap fonts.  You
  should use the equivalent AMS Fonts:
\begin{verbatim}
   \usepackage{amsfonts}
\end{verbatim}
followed by, e.g., \verb+\mathbb{R}+, \verb+\mathbb{N}+, or
\verb+\mathbb{C}+ for $\mathbb{R}$, $\mathbb{N}$ or $\mathbb{C}$.  You
can also use the following workaround for reals, natural and complex:
\begin{verbatim}
   \newcommand{\RR}{I\!\!R} %real numbers
   \newcommand{\Nat}{I\!\!N} %natural numbers
   \newcommand{\CC}{I\!\!\!\!C} %complex numbers
\end{verbatim}
Note that \verb+amsfonts+ is automatically loaded by the
\verb+amssymb+ package.

\end{itemize}

If your file contains type 3 fonts or non embedded TrueType fonts, we
will ask you to fix it.

\subsection{Margins in \LaTeX{}}

Most of the margin problems come from figures positioned by hand using
\verb+\special+ or other commands. We suggest using the command
\verb+\includegraphics+ from the \verb+graphicx+ package. Always
specify the figure width as a multiple of the line width as in the
example below:
\begin{verbatim}
   \usepackage[pdftex]{graphicx} ...
   \includegraphics[width=0.8\linewidth]{myfile.pdf}
\end{verbatim}
See Section 4.4 in the graphics bundle documentation
(\url{http://mirrors.ctan.org/macros/latex/required/graphics/grfguide.pdf})

A number of width problems arise when \LaTeX{} cannot properly
hyphenate a line. Please give LaTeX hyphenation hints using the
\verb+\-+ command when necessary.

\subsubsection*{Acknowledgments}

Use unnumbered third level headings for the acknowledgments. All
acknowledgments go at the end of the paper. Do not include
acknowledgments in the anonymized submission, only in the final paper.

\section*{References}

References follow the acknowledgments. Use unnumbered first-level
heading for the references. Any choice of citation style is acceptable
as long as you are consistent. It is permissible to reduce the font
size to \verb+small+ (9 point) when listing the references. {\bf
  Remember that you can use more than eight pages as long as the
  additional pages contain \emph{only} cited references.}
\medskip

\small

[1] Alexander, J.A.\ \& Mozer, M.C.\ (1995) Template-based algorithms
for connectionist rule extraction. In G.\ Tesauro, D.S.\ Touretzky and
T.K.\ Leen (eds.), {\it Advances in Neural Information Processing
  Systems 7}, pp.\ 609--616. Cambridge, MA: MIT Press.

[2] Bower, J.M.\ \& Beeman, D.\ (1995) {\it The Book of GENESIS:
  Exploring Realistic Neural Models with the GEneral NEural SImulation
  System.}  New York: TELOS/Springer--Verlag.

[3] Hasselmo, M.E., Schnell, E.\ \& Barkai, E.\ (1995) Dynamics of
learning and recall at excitatory recurrent synapses and cholinergic
modulation in rat hippocampal region CA3. {\it Journal of
  Neuroscience} {\bf 15}(7):5249-5262.

\end{document}
